\secnumbersection{CONCLUSIONES}
\label{sec:cap6}

Esta memoria evaluó el modelo \textit{Deep Gravity} para la distribución de flujos de movilidad en el transporte público de Santiago de Chile, con el objetivo de estudiar la sensibilidad del desempeño predictivo frente a la discretización espacial. A partir de una cohorte común de 60.511.977 viajes canónicos DTPM, observados entre el 1 y el 29 de noviembre de 2024 en un dominio de 34 comunas, se construyeron tres representaciones espaciales ---Zona~777 (793 unidades), H3-r7 (223 celdas) y H3-r8 (1.210 celdas)--- y se entrenaron nueve modelos mensuales con tres inicializaciones por rama. La evaluación comparó \textit{Deep Gravity} frente a tres referencias ---uniforme, de atracción entrenada y de gravedad exponencial--- mediante el \textit{Common Part of Commuters} (CPC), liberando la prueba una sola vez después de congelar todos los \textit{checkpoints} seleccionados por validación.

Los resultados principales del experimento mensual se resumen a continuación. \textit{Deep Gravity} obtuvo un CPC medio de prueba de 0,5609 en Zona~777, 0,7722 en H3-r7 y 0,6079 en H3-r8, con desviaciones estándar muestrales inferiores a 0,007 entre las tres inicializaciones. Estas cifras superan consistentemente a la referencia uniforme (0,2992, 0,3702 y 0,2611) y a la referencia de gravedad exponencial (0,3819, 0,6153 y 0,3777). Sin embargo, la referencia de atracción entrenada obtuvo un CPC mayor que la media de \textit{Deep Gravity} en las tres representaciones: 0,6574 frente a 0,5609 en Zona~777, 0,7757 frente a 0,7722 en H3-r7 y 0,6867 frente a 0,6079 en H3-r8. La diferencia fue particularmente estrecha en H3-r7, donde una inicialización individual de la red alcanzó 0,7787, superando el valor de la atracción; no obstante, el resumen de las tres inicializaciones no autoriza declarar superioridad del modelo sobre esta referencia.

La sensibilidad al peso de evaluación mostró que el cambio de masa expandida a conteos sin expansión no alteró la lectura comparativa. La conservación de masa se verificó con un error máximo absoluto por origen de $1{,}863 \times 10^{-9}$. En cuanto al análisis complementario de explicabilidad, las tres aproximaciones evaluadas ---gradientes integrados, SHAP por permutaciones agrupadas y valores de Shapley exactos para grupos jerárquicos--- no acreditaron todos los controles de completitud y estabilidad predefinidos, por lo que no se publican atribuciones, rankings ni figuras SHAP.

Un hallazgo transversal a los resultados es que la discretización espacial modifica el problema evaluado. Las tres representaciones comparten la misma cohorte de viajes, pero difieren en el número de destinos por origen, la densidad de pares OD positivos, la masa de prueba y la agregación de los atributos territoriales. El CPC no permite ordenar las representaciones como si fueran mediciones de una misma tarea, por lo que la comparación sustantiva se realiza dentro de cada rama. Esta dependencia del CPC respecto a la representación es en sí misma un resultado informativo para la comunidad de modelación de flujos.


\subsection{Cumplimiento de objetivos}

A continuación se examina el cumplimiento de cada uno de los seis objetivos específicos formulados en el Capítulo~1.

\textbf{Objetivo~1: Auditar y adaptar el código de referencia.} Se completó una auditoría técnica de la implementación de \textit{scikit-mobility} y se verificó su funcionamiento mediante una corrida de compatibilidad con datos de Nueva York. Esta corrida constituye un antecedente de compatibilidad de implementación, no una demostración de desempeño científico para Santiago. El pipeline resultante registra configuración efectiva, semillas, versiones, hashes de entradas, particiones, métricas y \textit{checkpoints} en un directorio exclusivo por corrida, lo que permite verificar la procedencia de cada resultado.

\textbf{Objetivo~2: Reconstruir viajes canónicos, definir una cohorte común y delimitar el núcleo urbano.} Se reconstruyeron los viajes canónicos mediante la unión de las tablas DTPM de viajes y etapas, exigiendo secuencia completa de etapas, bajada válida en la etapa final y coordenadas dentro de los rangos UTM de control. El dominio de estudio se delimitó a 34 comunas, con exclusión posterior de las zonas 848, 849 y 852 por carecer de polígono oficial verificable. La cohorte comparativa resultante contiene 60.511.977 viajes con 85.774.522,4070 de masa expandida, compartida por las tres ramas.

\textbf{Objetivo~3: Construir representaciones comparables y seleccionar candidatas.} Se evaluó la factibilidad de las resoluciones H3-r3 a H3-r8. Las resoluciones H3-r7 y H3-r8 pasaron al piloto junto con Zona~777 y se mantuvieron en el experimento mensual. Las tres representaciones operan sobre la misma cohorte, con adaptadores espaciales que asignan cada viaje a su unidad correspondiente sin alterar la cohorte fuente.

\textbf{Objetivo~4: Construir atributos con fuentes, unidades y procedimientos explícitos.} Se construyó un vector de 19 atributos por ubicación y 39 entradas por par OD, que incluye densidad de población censal (Censo 2024, interpolada por fracción de área), cinco coberturas de uso de suelo, tres densidades de longitud vial, cinco familias de servicios y la distancia geodésica. La normalización se estimó exclusivamente con orígenes activos de entrenamiento de cada rama y escenario, impidiendo que el conjunto de prueba interviniera en la estandarización.

\textbf{Objetivo~5: Entrenar y evaluar \textit{Deep Gravity} frente a tres referencias.} Se ejecutaron nueve corridas mensuales con tres representaciones y tres inicializaciones, utilizando el soporte completo de destinos. \textit{Deep Gravity} superó la referencia uniforme y la de gravedad exponencial en las tres ramas, con mejoras relativas del 87--109\% y del 25--47\% respectivamente. No obstante, la referencia de atracción entrenada obtuvo un CPC medio superior en las tres representaciones. Este resultado indica que, bajo la configuración evaluada, la distribución observada de llegadas a cada destino es una señal predictiva más fuerte que la combinación de atributos territoriales y distancia que utiliza la red. La referencia de atracción, sin embargo, requiere haber observado flujos previos para calcular la masa de llegada por destino, lo cual es precisamente la información que \textit{Deep Gravity} busca estimar a partir de atributos territoriales. Por ello, la comparación no invalida la utilidad del modelo en contextos donde no se dispone de datos de flujos previos, sino que delimita su ventaja comparativa.

\textbf{Objetivo~6: Evaluar la viabilidad de obtener atribuciones predictivas estables.} Se implementaron y evaluaron tres aproximaciones de XAI con controles predefinidos de completitud, convergencia y estabilidad. Los gradientes integrados no acreditaron completitud para una consulta crítica. El método por permutaciones no acreditó estabilidad entre fondos independientes en dos de tres inicializaciones. El método exacto jerárquico acreditó una primera consulta pero no una segunda consulta crítica. En consecuencia, este objetivo se cumplió parcialmente: se evaluó la viabilidad con rigor, pero el resultado es que no se obtuvieron atribuciones suficientemente estables para una interpretación sustantiva. No publicar atribuciones que no superan los controles es un resultado metodológico legítimo, no un fracaso del análisis.


\subsection{Contribuciones}

Las contribuciones de esta memoria se enumeran a continuación.

\begin{enumerate}
    \item \textbf{Primera evaluación reproducible de \textit{Deep Gravity} en una metrópolis latinoamericana.} Hasta donde se ha identificado en la literatura, este trabajo constituye la primera aplicación documentada de \textit{Deep Gravity} al contexto de Santiago de Chile, utilizando datos reales de transporte público (DTPM) y atributos territoriales de Censo 2024 y OpenStreetMap. La evaluación se realizó sobre tres representaciones espaciales y con un protocolo de prueba congelado.

    \item \textbf{Metodología de comparación entre discretizaciones con cohorte común.} El diseño experimental separa el núcleo común de viajes de los adaptadores espaciales, permitiendo estudiar la sensibilidad del desempeño predictivo a la discretización sin confundir cambios de datos con cambios de representación. Esta metodología es transferible a otras ciudades y modelos.

    \item \textbf{Pipeline con trazabilidad integral.} Cada corrida conserva semillas, hashes de entradas, versiones de código, particiones deterministas, métricas por origen y \textit{checkpoints} en un directorio exclusivo. La prueba se libera una sola vez después de congelar las decisiones de validación, sin selección post-hoc de semillas ni de representaciones. La conservación de masa se verifica explícitamente.

    \item \textbf{Resultado metodológico sobre los límites de XAI.} La evaluación sistemática de tres aproximaciones de explicabilidad, con controles predefinidos de completitud y estabilidad, aporta evidencia empírica de que obtener atribuciones estables en modelos de movilidad no es automático. Este resultado contribuye a la discusión sobre estándares de publicación de atribuciones en el área.

    \item \textbf{Caracterización de la referencia de atracción como \textit{baseline} dominante.} La identificación de que la distribución de atracción de destinos constituye la señal predictiva más fuerte en las tres representaciones es información útil para futuros trabajos que busquen superar este umbral mediante la incorporación de variables adicionales o arquitecturas alternativas.
\end{enumerate}


\subsection{Limitaciones}

Las limitaciones de este trabajo se agrupan en tres categorías.

\textbf{Limitaciones de datos.} La fuente de movilidad se restringe a los viajes reconstruibles desde DTPM durante las 29 fechas disponibles de noviembre de 2024. No cubre viajes en modos no observables (caminata pura, taxi, vehículo particular, bicicleta), por lo que los resultados no representan la totalidad de la movilidad urbana. El modelo requiere como entrada el flujo saliente observado por origen ($O_i$), lo cual limita su autonomía respecto a datos de validación. La cobertura de OpenStreetMap es heterogénea en Santiago, con mayor completitud en comunas de ingresos medios y altos, lo que puede sesgar los atributos de entrada en zonas periféricas. La instantánea OSM utilizada (2024-01-01) antecede a los viajes, y la interpolación areal del Censo 2024 supone distribución uniforme dentro de cada geometría censal.

\textbf{Limitaciones del diseño experimental.} Las tres inicializaciones caracterizan la variación de entrenamiento bajo una sola partición espacial fija (un \textit{tile} de prueba), por lo que no estiman incertidumbre poblacional ni variación entre particiones. No se dispone de una validación temporal independiente: el experimento mensual contiene el día del piloto y comparte su geografía de prueba. No se comparó \textit{Deep Gravity} contra modelos de oportunidades intervinientes (PWO, modelo unificado de Liu y Yan) ni contra arquitecturas de grafos (GNN), que el trabajo relacionado identifica como alternativas relevantes. El \textit{dropout} efectivo fue cero y no se exploró la sensibilidad a hiperparámetros de la arquitectura, como profundidad, ancho de capas o función de optimización.

\textbf{Limitaciones de alcance.} La evaluación distribuye flujos condicionados al flujo saliente observado; no predice la generación total de viajes por origen ni fechas futuras. Las conclusiones se limitan al transporte público observable en DTPM dentro de las 34 comunas del dominio de estudio. No se integró la generación de flujos con un proceso de planificación de transporte ni de optimización de rutas o frecuencias.


\subsection{Trabajo futuro}

Las líneas de trabajo futuro que se derivan de las limitaciones y hallazgos de esta memoria son las siguientes.

\textbf{Particiones espaciales múltiples.} Implementar una validación cruzada sobre \textit{tiles} (por ejemplo, \textit{k-fold} con $k$ igual al número de \textit{tiles}) permitiría estimar la incertidumbre espacial del CPC, no solo la variación asociada a la inicialización. Esto proporcionaría intervalos de confianza más robustos y una medida de la sensibilidad de los resultados a la elección del \textit{tile} de prueba.

\textbf{Validación temporal.} Entrenar con un período y evaluar con otro ---por ejemplo, entrenar con octubre y evaluar con noviembre, o utilizar entregas DTPM de otros meses--- permitiría evaluar la estabilidad del modelo ante cambios estacionales o variaciones post-pandemia en los patrones de movilidad, dimensión que esta memoria no cubre.

\textbf{Exploración de hiperparámetros.} La arquitectura utilizada hereda la configuración publicada sin explorar alternativas. Evaluar tasas de \textit{dropout} mayores que cero, profundidades y anchos de capa diferentes, tamaños de lote, y optimizadores alternativos (Adam frente a RMSprop) podría reducir la brecha con la referencia de atracción. El uso de GPU aceleraría esta exploración.

\textbf{Incorporación de modelos de oportunidades como referencias adicionales.} El modelo de oportunidades ponderadas por población (PWO) de Yan et al. \cite{yan2014population} y el modelo unificado de oportunidades de Liu y Yan \cite{liu2020intervening}, ambos libres de parámetros y diseñados para la escala intraurbana, constituyen referencias relevantes que no fueron evaluadas en esta memoria. Su inclusión permitiría situar el desempeño de \textit{Deep Gravity} dentro de una familia más amplia de modelos de flujos.

\textbf{Datos complementarios a OpenStreetMap.} La heterogeneidad de cobertura de OSM en Santiago podría mitigarse complementando los atributos con datos catastrales del Servicio de Impuestos Internos, datos agregados de telefonía móvil o registros agregados de la tarjeta bip!. Estas fuentes podrían mejorar la representación de las zonas periféricas donde OSM tiene menor cobertura.

\textbf{Estimación de la generación de viajes sin datos de flujo.} El modelo actual requiere el flujo saliente observado ($O_i$) como entrada. Una extensión natural es estimar $O_i$ a partir de datos de Censo, infraestructura de transporte y características territoriales, lo que haría al modelo más autónomo respecto a registros de validación previos y permitiría su aplicación en zonas o ciudades sin datos DTPM.

\textbf{Protocolos de explicabilidad más amplios.} El resultado negativo del análisis XAI sugiere que la combinación de una red profunda de 15 capas con un soporte amplio de destinos dificulta la obtención de atribuciones estables. Futuros trabajos podrían evaluar modelos más simples o más regularizados, aumentar el número de fondos de referencia, o emplear métodos de atribución que no requieran una referencia explícita (como los valores SHAP basados en \textit{kernel}).

\textbf{Integración con planificación de transporte.} Utilizar los flujos generados por \textit{Deep Gravity} como insumo para un modelo de optimización multiobjetivo de rutas ---por ejemplo, mediante NSGA-II--- permitiría evaluar la utilidad práctica de la matriz OD estimada en un contexto de planificación. Un análisis de sensibilidad que mida la variación en la calidad de las rutas frente a errores en la estimación de flujos completaría la cadena de valor del modelo.


\subsection{Reflexión final}

Esta memoria demuestra que \textit{Deep Gravity} puede ejecutarse de forma trazable y reproducible sobre datos reales de transporte público en Santiago de Chile, produciendo distribuciones de destino que superan las referencias uniforme y gravitacional en las tres discretizaciones evaluadas. El resultado de que la referencia de atracción entrenada domine en promedio no invalida el enfoque; por el contrario, delimita con precisión su ventaja comparativa: \textit{Deep Gravity} es valioso justamente en los contextos donde no se dispone de datos previos de flujo para construir una distribución de atracción, como en zonas de nuevo desarrollo urbano, ciudades sin encuestas de movilidad o escenarios de cambio estructural en los patrones de desplazamiento.

La sensibilidad del desempeño al nivel de discretización espacial es un recordatorio de que, en la modelación de flujos de movilidad, la representación del espacio no es un detalle técnico sino una decisión que condiciona el problema evaluado, los atributos que se construyen y la interpretación de los resultados. La cercanía observada en H3-r7, donde \textit{Deep Gravity} se aproxima a la atracción entrenada con una diferencia de 0,0035 en CPC, sugiere que con la resolución adecuada y con atributos territoriales accesibles, la red puede reconstruir una proporción sustancial de la estructura de destinos sin haber observado flujos reales hacia ellos.

El resultado negativo del análisis de explicabilidad es una contribución honesta que invita a la comunidad a definir estándares más explícitos para la publicación de atribuciones en modelos de movilidad. No toda predicción útil admite una explicación estable bajo las aproximaciones disponibles, y documentar esa limitación es preferible a presentar atribuciones cuya robustez no ha sido acreditada.

La trazabilidad y reproducibilidad del pipeline, desde la reconstrucción de la cohorte hasta la liberación de la prueba, permiten que otros investigadores repliquen, extiendan o refuten estos resultados con sus propios datos, particiones y configuraciones. Esta apertura metodológica es, en última instancia, la condición que permite que los hallazgos aquí presentados contribuyan al avance acumulativo del conocimiento en la modelación de flujos de movilidad urbana.
